\section{Conclusion}
\label{sec:conclusion}

This paper has examined whether an ECG front-end can determine in service, with measurements it
takes on itself, whether it still complies with IEC~60601-2-25, what has failed if it does not,
and whether the cause is in the circuit or in the electrodes. On a realistic design built around
an integrated instrumentation amplifier, a classifier on 26 self-test features reaches
\SI{0.9}{\percent} of escapes at \SI{7.8}{\percent} of false rejects, far from what a limit test
achieves, and three measurements taking three seconds come within one point of the balanced
accuracy of the complete self-test of \SI{97}{\second}. Defining a
fault by its effect on the specifications, as alternate test does, instead of by the deviation of
a component, as machine-learning diagnosis usually does, reduces false rejects from
\SI{37}{\percent} to about \SI{2}{\percent}. The cause is located with a macro F1 of 0.87 when
the answer is an ambiguity group obtained from the sensitivities of the nominal circuit, and that
answer, unlike the component, remains valid for fault magnitudes absent from training.
Non-compliant circuits are separated from electrode problems reliably; degraded electrode
contacts are not, because they fall within the variation between subjects.

The same analysis on a discrete three-op-amp front-end gives a worse picture, with structural
escapes and larger ambiguity groups, which shows that conclusions drawn on benchmark-like
circuits do not carry over to the circuits used in instruments.

The next step is validation on hardware: building the circuit, inducing a subset of the faults
and scoring the models trained on simulation against measured data. Multiple faults, a
per-patient reference for the electrode contact and the extension to other biopotential
front-ends are left for future work.
